\chapter{Conclusion and Summary}

This research confirms that while Fischer and Paterson established the theoretical foundations for string matching with wildcards, practical modern implementations heavily rely on Fast Fourier Transforms. Empirical verification with randomized trials throughout this research detected zero false positives, supporting the claim that the mathematical bounds established for Indyk's and Kalai's probabilistic methods hold consistently in practice.

To synthesize the theoretical complexities, empirical timings, and memory trade-offs evaluated in this research, a practical comparative overview for future software implementations is provided below:

\begin{itemize}
    \item \textbf{Naive:} This method is ideal for short patterns ($m < 32$) and small natural text alphabets. It requires zero heap allocation and offers excellent cache locality. However, its efficiency degrades significantly to $\mathcal{O}(nm)$ when processing dense wildcards or repetitive data sets.
    
    \item \textbf{FP Bitwise:} Serving primarily as a historical academic baseline, this method relies strictly on pure integer arithmetic. It is severely limited by the $\mathcal{O}(|\Sigma|^2)$ overhead generated by arbitrary-precision multiplication.
    
    \item \textbf{Baseline FFT:} Highly effective for data streams built on small alphabets (such as DNA, where $|\Sigma| \le 4$), delivering exact and fully deterministic matches. Its execution overhead is linearly dependent on the size of the underlying alphabet.
    
    \item \textbf{Sperner:} Designed for strictly deterministic applications requiring absolute certainty. It guarantees algorithmic performance mathematically bounded by a logarithmic alphabet scaling of $\mathcal{O}(\log |\Sigma|)$. Implementing it introduces mild combinatorial overhead needed for dynamic subset generation.
    
    \item \textbf{Indyk:} Suitable for traversing large character sets such as Unicode. It is asymptotically almost optimal, operating in $\mathcal{O}(n \log n)$ time, completely independent of the alphabet size. In practical software, however, it is burdened by a high constant factor generated by the need to compute independent FFT passes per validation.
    
    \item \textbf{Kalai:} A highly efficient architectural choice for practical applications requiring fast convolutions. It maintains a minimum constant factor of exactly 4 FFT iterations, ensuring execution speeds bounded by $\mathcal{O}(n \log m)$. It requires system infrastructure explicitly tolerant of Monte Carlo false positive rates capped at $\le 1/N$.
    
    \item \textbf{Clifford-Clifford:} The optimal choice for purely deterministic $\mathcal{O}(n \log m)$ pattern matching. By expanding Euclidean distance into three separate transforms, it provides absolute accuracy independent of the alphabet size, avoiding probabilistic errors at the cost of a slightly heavier constant factor than Kalai.
\end{itemize}

The empirical results definitively show that the Fast Fourier Transform fundamentally restructures wildcard string matching, providing robust scalability for modern data environments.